\documentclass[11pt,a4paper]{article}
\usepackage[utf8]{inputenc}
\usepackage[T1]{fontenc}
\usepackage{lmodern}
\IfFileExists{french.ldf}{\usepackage[french]{babel}}{%
  \AtBeginDocument{\renewcommand{\abstractname}{Résumé}\renewcommand{\refname}{Références}\renewcommand{\tablename}{Tableau}}}
\usepackage[margin=2.2cm]{geometry}
\usepackage{booktabs,graphicx,hyperref,xcolor,listings,float}
\hypersetup{colorlinks=true,linkcolor=blue!50!black,urlcolor=blue!50!black}
\lstset{basicstyle=\ttfamily\small,breaklines=true,frame=single,columns=fullflexible}
% Valeur à compléter après l'entraînement (affichée en rouge tant qu'elle n'est pas remplie)
\newcommand{\TODO}[1]{\textcolor{red}{#1}}

\title{Détection d'équipements de protection individuelle en temps réel\\
\large Rapport technique — Module Deep Learning Avancé (Axe 3, YOLO)}
\author{Omar Ghannem\\
Cycle ingénieur IA \& Data Science, 3\textsuperscript{e} année — TEK-UP\\
Enseignant : Haythem Ghazouani}
\date{\today}

\begin{document}
\maketitle
\begin{abstract}
Nous présentons un système complet de détection de personnes et d'équipements de protection
individuelle (EPI) en milieu de travail, de l'entraînement au déploiement web. Des modèles YOLO
sont entraînés sur le dataset public SH17 (8\,099 images, 17 classes), suivis avec MLflow,
exportés en ONNX et servis par une API FastAPI consommée par une interface React avec webcam.
Une règle de conformité déclenche une alerte lorsqu'une tête est détectée sans casque.
La chaîne est reproductible (DVC, MLflow, Docker Compose) et testée par GitHub Actions.
Le meilleur modèle (YOLOv8n à 960\,px) atteint 56,9\,\% de mAP@50 et 35,0\,\% de mAP@[.5:.95] sur le jeu de test, proche des 58,0\,\% publiés, avec 20 epochs seulement ; l'export ONNX divise la latence CPU par 1,5.
\end{abstract}

\section{Problématique et objectifs}
Le contrôle du port des EPI est aujourd'hui surtout manuel. Problématique : \emph{comment détecter
automatiquement et en temps réel, via une application web, les personnes et leurs équipements de
protection, et alerter en cas de non-conformité ?}

Objectifs mesurables : (1) baseline YOLO nano proche des 58,0\,\% mAP@50 publiés pour YOLOv8n ;
(2) comparaison d'au moins deux architectures sur mAP@[.5:.95] et latence ; (3) effet de la
résolution d'entrée (640 contre 960) ; (4) export ONNX et gain de latence mesuré ; (5) API et
frontend avec webcam et alerte.

\section{Données}
\subsection{Le dataset SH17}
SH17 \cite{sh17} contient 8\,099 images issues de Pexels, 75\,994 instances et 17 classes, au format
YOLO, sous licence CC BY-NC-SA 4.0 (usage académique). L'ordre des identifiants est celui du fichier
\texttt{sh17.yaml} (alphabétique), différent du README ; nous l'avons vérifié visuellement.

\subsection{Exploration et qualité}
\begin{itemize}
  \item Aucune image corrompue, aucune annotation invalide, aucun doublon, aucune fuite train/val.
  \item Fort déséquilibre : ratio max/min de 118,3 (\texttt{hands} 15\,850 contre \texttt{face-guard} 134) ;
        les EPI clés sont rares (\texttt{helmet} 927, \texttt{safety-vest} 530).
  \item Petits objets : 35,7\,\% des boîtes font moins de 32\,px à 640\,px d'entrée (oreilles 92\,\%, pieds 72\,\%).
  \item Densité moyenne de 9,4 objets par image ; résolutions de 1335$\times$1002 à 8192$\times$8192.
  \item Contenu : surtout des bureaux et ateliers, peu de chantiers.
\end{itemize}

\subsection{Protocole de découpage}
\begin{table}[H]\centering
\begin{tabular}{llr}\toprule
Split & Origine & Images\\\midrule
train & 90\,\% du train officiel & 5\,831\\
val & 10\,\% du train officiel (seed 42) & 648\\
test & val officiel (comparable au benchmark) & 1\,620\\\bottomrule
\end{tabular}\end{table}
Le découpage est fait une seule fois, avant toute augmentation, par le stage DVC \texttt{prepare}
qui réduit aussi les images à 1\,280\,px de grand côté (les labels normalisés restent valides).
Le choix du meilleur modèle se fait sur la validation ; le test ne sert qu'au rapport final.

\section{Modélisation}
Nous utilisons la famille YOLO (Ultralytics) en transfer learning depuis COCO : un réseau de
détection en une passe, adapté au temps réel. Trois runs, même seed et même budget :
\begin{enumerate}
  \item \textbf{baseline} YOLOv8n, 640\,px, 30 epochs (même modèle que le benchmark publié) ;
  \item \textbf{YOLO26n}, 640\,px, 30 epochs : architecture récente, sortie sans NMS ;
  \item YOLOv8n à \textbf{960\,px}, 20 epochs : effet de la résolution sur les petits objets.
\end{enumerate}
Augmentations Ultralytics par défaut (mosaïque, HSV, flips, translation/échelle), early stopping
(patience 10), sélection sur le mAP@[.5:.95] de validation.

\section{Résultats}
\begin{table}[H]\centering
\begin{tabular}{lrrrrr}\toprule
Run & val mAP50-95 & test mAP50 & test mAP50-95 & AP50 helmet & GPU (ms)\\
 & (\%) & (\%) & (\%) & (\%) & \\\midrule
baseline-yolov8n-640 & 32,0 & 52,7 & 32,4 & 60,7 & 3,1\\
yolo26n-640 & 32,3 & 53,3 & 33,2 & 62,3 & 3,1\\
\textbf{yolov8n-960} & \textbf{33,7} & \textbf{56,9} & \textbf{35,0} & \textbf{68,5} & 7,5\\
\midrule
Référence YOLOv8n \cite{sh17} & -- & 58,0 & 36,6 & -- & --\\\bottomrule
\end{tabular}
\caption{Comparaison des modèles (source : \texttt{reports/comparison.csv}, MLflow).}
\end{table}

\paragraph{Lecture du tableau.} Le choix se fait sur la validation : \texttt{yolov8n-960} est le meilleur
(33,7\,\% de mAP@[.5:.95]) et le test le confirme. À même résolution, YOLO26n fait un peu mieux que
YOLOv8n (+0,6 point de mAP@50, +0,8 point de mAP@[.5:.95]) pour la même latence GPU. Le passage de 640 à
960\,px apporte le gain le plus net (+4,2 points de mAP@50, +7,8 points d'AP@50 sur \texttt{helmet}), au prix
d'une latence GPU 2,5 fois plus élevée. À 640\,px, notre baseline (52,7\,\%) reste sous les 58,0\,\% publiés :
nous n'avons fait que 30 epochs.

\paragraph{Analyse par taille d'objet.} Rappel à IoU 0,5 du meilleur modèle sur le test (taille vue par le
réseau, racine de l'aire de la boîte) :

\begin{table}[H]\centering
\begin{tabular}{lrrr}\toprule
Taille & Objets & Détectés & Rappel\\\midrule
$<$16\,px & 1\,472 & 337 & 22,9\,\%\\
16--32\,px & 2\,103 & 1\,324 & 63,0\,\%\\
32--96\,px & 5\,518 & 4\,354 & 78,9\,\%\\
$>$96\,px & 6\,265 & 5\,501 & 87,8\,\%\\\bottomrule
\end{tabular}\end{table}
Les très petits objets sont la principale source d'erreurs : moins d'un sur quatre est retrouvé, contre près
de 9 sur 10 pour les grands. C'est ce qui explique le gain apporté par la résolution 960\,px.

\paragraph{Analyse par classe.} Les EPI clés restent les plus difficiles : \texttt{helmet} atteint 68,5\,\%
d'AP@50 mais \texttt{safety-vest} seulement 39,7\,\% (530 instances, la classe EPI la plus rare après
\texttt{face-guard} et \texttt{medical-suit}). La règle d'alerte reposant sur \texttt{helmet}, un casque
manqué produit une fausse alerte.

\section{Export ONNX et latence}
Le meilleur modèle (\texttt{yolov8n-960}) est exporté en ONNX (opset 17, entrée fixe 960$\times$960).
Sur le CPU de Kaggle, à entrée identique : PyTorch 216,9\,ms contre ONNX Runtime 141,6\,ms (médiane sur 50
passages après 10 d'échauffement), soit un gain de 1,53$\times$. L'écart de mAP@50 entre ONNX et PyTorch sur le
test étant inférieur à 2 points, la version 2 de \texttt{sh17-ppe-detector} est passée de l'alias
\texttt{staging} à \texttt{production} dans le MLflow Model Registry (la version 1, YOLO26n, issue du premier
export, reste disponible). Sur un ordinateur portable, l'API mesure environ 190\,ms d'inférence à 960\,px,
soit environ 4 images par seconde en webcam ; le modèle à 640\,px serait environ 2 fois plus rapide.

\section{Architecture logicielle et MLOps}
\begin{itemize}
  \item \textbf{DVC} : \texttt{dvc.yaml} (stage \texttt{prepare}), \texttt{params.yaml}, \texttt{dvc.lock}
        (empreintes des données) ; aucune donnée ni aucun poids dans Git.
  \item \textbf{MLflow} : paramètres, courbes par epoch, métriques de test par classe, artefacts
        (poids, matrice de confusion, analyse par taille), Model Registry avec alias.
  \item \textbf{API FastAPI} : \texttt{/health}, \texttt{/model/metadata}, \texttt{/predict},
        \texttt{/predict/visualize} ; schémas Pydantic, Swagger sur \texttt{/docs}, erreurs 400 et 503,
        modèle ONNX chargé une seule fois au démarrage, inférence hors de la boucle asynchrone.
  \item \textbf{Frontend React} : glisser-déposer, galerie d'exemples, webcam, canvas de boîtes
        colorées par classe, bandeau d'alerte, latences client et serveur.
  \item \textbf{Docker Compose} : services \texttt{mlflow}, \texttt{backend}, \texttt{frontend}
        (images multi-stage) lancés par \texttt{docker-compose up --build}.
  \item \textbf{GitHub Actions} : black, flake8, 18 tests pytest, eslint, build du frontend, build Docker.
\end{itemize}

\paragraph{Règle d'alerte.} Une tête est conforme si le centre d'un casque se trouve dans sa boîte
agrandie de 50\,\% vers le haut. On part de \texttt{head} (abondante) plutôt que de \texttt{person}.

\section{Limites et perspectives}
\begin{itemize}
  \item Classes EPI rares (\texttt{helmet}, \texttt{safety-vest}) : fausses alertes possibles ;
        piste : sur-échantillonnage ciblé ou ajout d'un dataset de chantiers.
  \item Budget GPU limité (Kaggle) : 20 à 30 epochs seulement, loin d'un entraînement complet.
  \item Scènes surtout de bureaux : la démo webcam détecte bien \texttt{head}, \texttt{face}, \texttt{hands},
        moins bien les casques sans accessoire réel.
  \item Perspectives : modèle \texttt{s}, résolution 1280, suivi multi-objets, quantification INT8.
\end{itemize}

\begin{thebibliography}{9}
\bibitem{sh17} H.~M. Ahmad, A.~Rahimi, \emph{SH17: A Dataset for Human Safety and Personal Protective
Equipment Detection in Manufacturing Industry}, arXiv:2407.04590, 2024.
\bibitem{yolo} G.~Jocher et al., \emph{Ultralytics YOLO}, \url{https://github.com/ultralytics/ultralytics}.
\end{thebibliography}
\end{document}
